\section{Related Work}
It has been shown that global features or contextual interactions \cite{he2004multiscale,shotton2009textonboost,kohli2009robust,ladicky2009associative,gould2009decomposing,yao2012describing} are beneficial in correctly classifying pixels for semantic segmentation. In this work, we discuss four types of Fully Convolutional Networks (FCNs) \cite{sermanet2013overfeat, long2014fully} (see \figref{fig:arch} for illustration) that exploit context information for semantic segmentation \cite{hariharan2014hypercolumns, dai2014convolutional, mostajabi2014feedforward, chen2015semantic, zheng2015conditional, lin2015efficient, papandreou2015weakly, schwing2015fully, xia2015zoom}.

\textbf{Image pyramid:} The same model, typically with shared weights, is applied to multi-scale inputs. Feature responses from the small scale inputs encode the long-range context, while the large scale inputs preserve the small object details. Typical examples include Farabet \etal \cite{farabet2013learning} who transform the input image through a Laplacian pyramid, feed each scale input to a DCNN and merge the feature maps from all the scales. \cite{eigen2014predicting, pinheiro2014recurrent} apply multi-scale inputs sequentially from coarse-to-fine, while \cite{lin2015efficient,chen2015attention,chen2016deeplab} directly resize the input for several scales and fuse the features from all the scales. The main drawback of this type of models is that it does not scale well for larger/deeper DCNNs (\eg, networks like \cite{he2015deep,zagoruyko2016wide,wu2016wider}) due to limited GPU memory and thus it is usually applied during the inference stage \cite{dai2015boxsup}.

\textbf{Encoder-decoder:} This model consists of two parts: (a) the encoder where the spatial dimension of feature maps is gradually reduced and thus longer range information is more easily captured in the deeper encoder output, and (b) the decoder where object details and spatial dimension are gradually recovered. For example, \cite{long2014fully,noh2015learning} employ deconvolution \cite{zeiler2011adaptive} to learn the upsampling of low resolution feature responses. SegNet \cite{badrinarayanan2015segnet} reuses the pooling indices from the encoder and learn extra convolutional layers to densify the feature responses, while U-Net \cite{ronneberger2015u} adds skip connections from the encoder features to the corresponding decoder activations, and \cite{ghiasi2016laplacian} employs a Laplacian pyramid reconstruction network. More recently, RefineNet \cite{lin2016refinenet} and \cite{pohlen2016full, peng2017large, islamgated} have demonstrated the effectiveness of models based on encoder-decoder structure on several semantic segmentation benchmarks. This type of model is also explored in the context of object detection \cite{lin2016feature,shrivastava2016beyond}.

\textbf{Context module:} This model contains extra modules laid out in cascade to encode long-range context. One effective method is to incorporate DenseCRF \cite{krahenbuhl2011efficient} (with efficient high-dimensional filtering algorithms \cite{adams2010fast}) to DCNNs \cite{chen2014semantic, chen2016deeplab}. Furthermore, \cite{zheng2015conditional, lin2015efficient, schwing2015fully} propose to jointly train both the CRF and DCNN components, while \cite{liu2015semantic, yu2015multi} employ several extra convolutional layers on top of the belief maps of DCNNs (belief maps are the final DCNN feature maps that contain output channels equal to the number of predicted classes) to capture context information. Recently, \cite{jampani2016learning} proposes to learn a general and sparse high-dimensional convolution (bilateral convolution), and \cite{Vemulapalli2016Gaussian, chandra2016fast} combine Gaussian Conditional Random Fields and DCNNs for semantic segmentation.

\textbf{Spatial pyramid pooling:} This model employs spatial pyramid pooling \cite{grauman2005pyramid, lazebnik2006beyond} to capture context at several ranges. The image-level features are exploited in ParseNet \cite{liu2015parsenet} for global context information. DeepLabv2 \cite{chen2016deeplab} proposes atrous spatial pyramid pooling (ASPP), where parallel atrous convolution layers with different rates capture multi-scale information. Recently, Pyramid Scene Parsing Net (PSP) \cite{zhao2016pyramid} performs spatial pooling at several grid scales and demonstrates outstanding performance on several semantic segmentation benchmarks. There are other methods based on LSTM \cite{hochreiter1997long} to aggregate global context \cite{liang2015semantic, byeon2015scene, yan2016combining}. Spatial pyramid pooling has also been applied in object detection \cite{he2014spatial}.

In this work, we mainly explore atrous convolution \cite{holschneider1989real, giusti2013fast, sermanet2013overfeat, papandreou2014untangling, chen2014semantic, yu2015multi, chen2016deeplab} as a {\bf{context module}} and tool for {\bf{spatial pyramid pooling}}. Our proposed framework is general in the sense that it could be applied to any network. To be concrete, we duplicate several copies of the original last block in ResNet \cite{he2015deep} and arrange them in cascade, and also revisit the ASPP module \cite{chen2016deeplab} which contains several atrous convolutions in parallel. Note that our cascaded modules are applied directly on the feature maps instead of belief maps. For the proposed modules, we experimentally find it important to train with batch normalization \cite{ioffe2015batch}. To further capture global context, we propose to augment ASPP with image-level features, similar to \cite{liu2015parsenet, zhao2016pyramid}.

\textbf{Atrous convolution:} Models based on atrous convolution have been actively explored for semantic segmentation. For example, \cite{wu2016bridging} experiments with the effect of modifying atrous rates for capturing long-range information, \cite{wang2017understanding} adopts hybrid atrous rates within the last two blocks of ResNet, while \cite{dai2017deformable} further proposes to learn the deformable convolution which samples the input features with learned offset, generalizing atrous convolution. To further improve the segmentation model accuracy, \cite{wanglearning} exploits image captions, \cite{jain2017fusionseg} utilizes video motion, and \cite{kong2017recurrent} incorporates depth information. Besides, atrous convolution has been applied to object detection by \cite{papandreou2014untangling, dai2016rfcn, huang2016speed}.
